% Supporting Information for EDGE paper 3 (Statistics in Medicine).
% The material moved out of the main text in the restructure of 2026-09-30, unchanged in substance.
% Build: pdflatex supplement -> bibtex supplement -> pdflatex supplement x2
\documentclass[11pt]{article}
\usepackage[a4paper,margin=2.5cm]{geometry}
\usepackage{amsmath,amssymb}
\usepackage{booktabs,multirow,adjustbox}
\usepackage[table]{xcolor}
\usepackage[round]{natbib}
% the natbib.sty in this folder is Wiley's copy, whose \bibsection calls the USG class's \bmsection
\renewcommand{\bibsection}{\section*{References}}
\makeatletter
\def\input@path{{tables/}}
\makeatother
\renewcommand{\thetable}{S\arabic{table}}
\renewcommand{\thesection}{S\arabic{section}}
\setcounter{section}{-1}
\newcommand{\EFname}{\textsf{HL}$_F$}
\newcommand{\pibar}{\overline{\hat{\pi}}}
\providecommand{\qed}{\ensuremath{\square}}

\title{Supporting Information for\\ ``A grouped calibration test for logistic regression that tolerates a few corrupted records''}
\author{Ebrahim Khaled Ebrahim, Osama Abd El-Aziz Hussein and Ahmed El-Kotory}
\date{}

\begin{document}
\maketitle

\section{The comparisons specified before the study was run}\label{si:hyp}

Before any scenario was computed, seven comparisons were written down with their scenarios, their
statistic, a threshold and a rule for what would count against EDGE, and the document
was fixed by its SHA-256 hash. Table~\ref{tab:hypotheses} lists them with their outcomes. Three went
against the test: H1 and H3, where the GiViTI belt is the better test when events are scarce or the
curve bends once, and H4, where EDGE at ten groups led or tied the best of four partition
tests in $17$ of the $19$ selected scenarios, where all $19$ were required; at the default partition it
did so in all $21$. The main text reports each of them as a finding in Section~5.2 and as
a usage boundary in Section~8.3. The hash-fixed documents are in the reproducibility archive.

The main study re-ran, under these documents, a battery that had been run before them. The comparison
with robust tests (Section~6.6), the external-validation, streaming and what-protects studies were
added after the main results had been read, each fixed by its own hash before any of its replicates
ran. Its rule was that the Mallows form of the robust Stukel test
would be said to match EDGE's protection if its false-alarm rate stayed at or below $0.10$
with ten corrupted records in $1000$ under both errors. It did not: $0.168$ and $0.789$.

\input{tab_hypotheses}

\section{The choice of weighting}\label{si:weighting}

Which risk groups carry the test depends on the weighting, and the right weighting depends on the
basis. For the symmetric basis the score weighting is what makes the test work at high
discrimination. It leads the unit form in $21$ of the $32$ power scenarios, and against a probit truth
the difference is large: power rises from $0.096$ to $0.598$ at $n=3500$, and from $0.506$ to $0.834$
at $n=5900$, at nominal size in both forms. These are the two scenarios in which this basis, in its
unit form, falls behind Stukel's joint score test by $0.352$ and $0.214$. With the score weighting the
order reverses and EDGE leads by $0.150$ and $0.114$.

Scaling by $\sqrt{V_g}$ moves weight towards the extreme risk groups, and the symmetric shape
$\eta|\eta|$ is carried by exactly those groups. The cubic basis already reaches them through its own
columns, so the same rescaling buys it nothing and costs it variance: for the cubic basis the score
form is behind in $24$ of the $32$ scenarios and ahead in one, by a mean of $0.039$.

One regime overrides both rules. The score weighting leans on the extreme groups, and those are the
groups a corrupted record occupies: with $25$ of $1000$ covariates corrupted, the cubic basis in its
score form reached a false-alarm rate of $0.894$ where its unit form reached $0.344$. When the
predictors may contain corrupted records, the unit form is required whatever the basis.

\section{The three further corruptions}\label{si:c2}

When the $k$ highest-risk records have their outcome set to zero (C2), every test rises together, 
EDGE included: at ten such records in a thousand it reads $0.667$ against Stukel's $0.718$
and the cubic test's $0.837$. A missed event at the top of the risk scale is the pattern a
calibration test is built to find.

A unit error that makes a prediction milder rather than extreme (C3) moves no test under the logistic
truth, as the mechanism of Section~3 of the main text predicts, since such a record lands in a crowded
middle group rather than a sparse extreme one. Under the two misfit truths it costs some power:
sixteen of the $168$ test-by-scenario readings move by more than three nominal standard errors, ten
under the complementary log--log truth and six under the probit, and EDGE is among them.
The largest move is its own, $0.818$ falling to $0.770$ at ten corrupted records under the
complementary log--log truth. These moves are losses of power against a real misfit, not false
alarms. Duplicated high-risk records (C4) leave every test between $0.054$ and $0.071$, and the
Hosmer--Lemeshow variants between $0.040$ and $0.057$.

\section{The resampling tests under corruption}\label{si:rivals}

The two resampling tests are too slow for the main contamination design, so a separate study put them
on a smaller one: three records in five hundred with their covariate multiplied by four or by eight,
$100$ replicates a scenario, and every test on the same data sets. With $100$ replicates the Monte
Carlo standard error of a rate of $0.05$ is $0.022$, so a rise of a few points cannot be seen. Multiplying by eight carries a
corrupted record's linear predictor to a median of $11$ and at most $16.5$.
Table~\ref{tab:rivals-contam} gives the result.

BAGofT and EDGE with the cubic basis gave no sign of moving: $0.06$, $0.01$ and $0.05$ for
the first in the clean, $\times4$ and $\times8$ scenarios, and $0.08$, $0.02$ and $0.05$ for the second;
with $100$ replicates a rise of a few points could not be seen. Le Cessie's smoothed test reads $0.07$, $0.04$ and $0.07$. The Liu projection test holds at
$\times4$ and breaks at $\times8$, where it reads $0.22$; on the same replicates it raised $18$ false
alarms that EDGE did not, against one the other way (Holm-corrected McNemar $p<0.001$).
The record-level score and calibration tests sit between $0.29$ and $0.49$ in both corrupted
scenarios. The symmetric basis reads $0.20$ at $\times8$. This is consistent with the main text: the
cubic basis, here at the default partition (twenty groups), did not move, while the symmetric basis
did at $\times8$.

\input{tab_rivals_contam}

\section{The robust fit}\label{si:robust}

Every data set of this study was fitted twice, by maximum likelihood and by the Huber-type
quasi-likelihood estimator of \citet{cantoni2001robust} (Huber's $\psi$ on the Pearson residual at
$1.345$, no weight on the covariates, \textsf{robustbase}'s default), which bounds the influence of the
residual but not of the covariates, and every
test was applied to each fit. Table~\ref{tab:robust} gives the result.

Five of the ten tests fail to hold their level after a robust fit on uncontaminated data: Stukel's two
score tests, at $0.199$ and $0.223$, and three of EDGE's four forms, the symmetric basis
in its unit form ($0.164$) and both score-weighted forms ($0.619$ and $0.225$). The rule fixed in
advance was that a failed level check leaves the study measuring that mismatch rather than the
contamination; we depart from it in one restricted way by reading the contaminated columns for the five
tests whose level survives. A robust fit does not bring the record-level tests back to the nominal
level.

Of those five, none has its false alarms reduced. At ten corrupted records the GiViTI belt reads $0.379$
under maximum likelihood and $0.400$ under the robust fit, the cubic calibration test $0.465$ and
$0.466$, and EDGE $0.089$ and $0.087$. Several are worse at twenty-five: the
Hosmer--Lemeshow statistic rises from $0.114$ to $0.171$ and the belt from $0.639$ to $0.678$.
A corrupted record's prediction is formed at its corrupted covariate, so refitting does not reach it;
restoring the slope, if anything, makes that prediction more extreme. Two further arms change nothing: at $n=5000$ with twenty-five corrupted
records EDGE reads $0.144$ under both fits, and against a weak probit departure the
robust fit does not pay for itself either. That a robust fit does not repair a non-robust test is the
observation on which robust testing was founded \citep{hampel1986robust,heritier1994robust}, and
\citet{cantoni2001robust} pair their estimator with robust quasi-deviance tests for that reason.

\input{tab_robust}

\section{Partitions that pool the extreme groups}\label{si:remedies}

Corrupted covariates produce extreme predictions, so they collect in the two extreme groups. A partition
that pools only those groups, and keeps the default partition in the middle where most of the power
lies, was studied as an alternative to ten groups: \textsf{EDGE-FR}$(c)$ puts the outer $5c$ of the
records at each end of the risk scale into one group. It does not have equal-frequency groups, so its
level was checked by simulation; it reads $0.039$ and $0.050$ on clean data at the two sample sizes,
and $0.041$ and $0.054$ on a fresh set of replicates. Table~\ref{tab:edgefr} gives the result. At
$c=1\%$ it matches ten groups up to ten corrupted records in $1000$ and fifty in $5000$ ($0.081$ against
$0.085$, and $0.111$ against $0.107$), is less protected beyond ($0.140$ against $0.101$ at twenty in
$1000$, $0.431$ against $0.328$ at a hundred in $5000$), and has the same power: a size-adjusted mean of
$0.490$ against $0.483$ over the six link and tail scenarios at $n=1000$, where the default partition has $0.500$.
A wider tail costs power, a mean of $0.449$ at $c=2\%$ and $0.408$ at $c=4\%$, without protection in
proportion: at $c=2\%$ it
reads $0.104$ with twenty corrupted records in $1000$ but $0.338$ with a hundred in $5000$, the same
rate. The main text therefore recommends the simpler ten-group partition.

\input{tab_edgefr}

Sharper remedies were considered first. Dropping the outer one or two per
cent of records, or dropping the outer groups outright, protects almost completely, between $0.047$
and $0.059$ at $n=5000$ with twenty-five and with fifty corrupted records, and costs almost all the
power, a mean raw rejection rate of $0.046$ to $0.059$ over the six link and tail scenarios against
the default's $0.508$. Pooling a tail of $1\%$ rather than $5\%$ keeps the power ($0.503$) and does not protect
($0.271$ at fifty corrupted records in $5000$). The extremes carry the signal, so a remedy that
removes them removes the departure with the corruption.

\section{False alarms of every test under the exaggerated covariate}\label{si:contamination}

Table~\ref{tab:contamination} gives, for every test of the main comparison, the false-alarm rate under
corruption C1 of the main text: $k$ records whose covariate is multiplied by four, in a sample for which
the model is otherwise correct. Section~6.1 of the main text reads it.

\input{tab_contamination}

\section{External validation against the usual tests}\label{si:external}

When a published model is checked on new patients, the usual tests are Cox's recalibration test,
Spiegelhalter's $z$ and the GiViTI belt in its external mode. Table~\ref{tab:external} compares 
EDGE in external mode with them, on a design fixed and hashed before it was run. All seven
tests hold their level. Over the ten departures EDGE at ten groups has a mean size-adjusted
power of $0.415$, against $0.420$ for the belt, $0.410$ for Cox's test, $0.335$ for the Hosmer--Lemeshow
statistic, $0.327$ for Stukel's test and $0.271$ for Spiegelhalter's $z$. The averages hide the shape of
the departure. Paired on the same replicates, EDGE at ten groups minus Cox's test is $-0.079$ and
$-0.042$ (Monte Carlo standard errors $0.011$ and $0.007$) on the two shifts, $-0.021$ and $-0.015$ on
the slopes, $-0.068$ and $-0.141$ on the threshold terms, $+0.045$ and $+0.118$ on the U-shaped terms
and $+0.053$ and $+0.206$ on the interactions, with standard errors between $0.004$ and $0.015$; the
differences from the belt are within $0.07$ of these. The mean difference is $+0.006$ against Cox's test
and $-0.005$ against the belt, and without the two interactions $-0.025$ and $-0.030$.

The specification wrote three expectations in advance. That Cox's test, Spiegelhalter's $z$ and the belt
would lead on calibration in the large and the slope held for Cox's test and the belt, and for
Spiegelhalter's $z$ only on the slope; on calibration in the large it is last ($0.107$ and $0.303$). That
EDGE would be competitive on the U-shaped and threshold terms held for the U-shape and failed for the
threshold, where it trails by up to $0.141$. That no test would detect the interaction well failed:
EDGE has a size-adjusted power of $0.459$ against it at the stronger setting at the default partition,
and $0.424$ at ten groups.

Under corrupted records EDGE at ten groups keeps every rate at or below $0.082$, below the bar of
$0.10$ though above the nominal level. Cox's test reads $0.088$ with five reversed records and $0.158$
with ten, the belt $0.192$ and $0.347$, and Stukel's test $0.610$ with ten. The Hosmer--Lemeshow
statistic and Spiegelhalter's $z$ stay below $0.10$ in every setting, with less power. Cox's test and the
belt are not two independent comparators: the belt fits a polynomial of the logit of the prediction
whose degree it selects by forward tests at its default threshold of $0.95$, and at degree one its model is Cox's
recalibration model, so in about $95\%$ of null replicates, and in $57\%$ to $96\%$ under the
alternatives, it returns Cox's $p$-value. Its extra false alarms arise when corrupted records push the
selected degree above one.

\input{tab_external}

\section{EDGE on streaming data}\label{si:stream}

\paragraph{Setting.} A frozen model gives each patient a prediction $p_i\in(0,1)$ before the outcome
$y_i\in\{0,1\}$ is seen. Cut points $0<c_1<\dots<c_{G-1}<1$ are fixed in advance, from reference
predictions or by the analyst, and do not depend on any outcome; patient $i$ falls in group
$g(i)=1+\#\{j:c_j<p_i\}$. After $n$ patients the monitor holds, for each group,
\[
O_g=\sum_{g(i)=g}y_i,\qquad E_g=\sum_{g(i)=g}p_i,\qquad V_g=\sum_{g(i)=g}p_i(1-p_i),\qquad m_g=\#\{i:g(i)=g\},
\]
and computes the external-mode statistic of Section~2 of the main text from them alone:
$r_g=(O_g-E_g)/\sqrt{V_g}$ and $\pibar_g=E_g/m_g$ over the non-empty groups, $Z$ the chosen basis evaluated
at $\pibar_g$ with a constant column, reduced to full column rank $k$, and $S_n=r^\top P_Z r$.

\paragraph{Proposition.} (a) $S_n$ is the same for every order and every division into batches in
which the $n$ patients are added, and equals the one-shot external statistic computed on all $n$
patients with the same groups. (b) Suppose that, given all the predictions, past and future, the outcomes
are independent with $\Pr(y_i=1)=p_i$ (the model is calibrated), that $G$ is fixed, and that
$V_g\to\infty$ in every group that is non-empty. Let $k_n$ be the rank of $Z$, which equals $d+1$ once
$d+1$ groups with distinct mean risks are non-empty. Then
$\sup_x|\Pr(S_n\le x\mid p_1,\dots,p_n)-F_{\chi^2_{k_n}}(x)|\to0$, for any sequence of predictions under
which every non-empty group's variance grows, however unequal the groups. (c) Adding a patient costs
$O(\log G)$ operations and computing $S_n$ costs $O(Gk^2)$, independently of $n$.

\paragraph{Proof.} (a) Each of $O_g,E_g,V_g,m_g$ is a sum over patients of a quantity that depends on
that patient alone, so it does not depend on the order of the terms or on how they are grouped, and
$S_n$ is a function of these sums only.

(b) Condition on the predictions, so that the groups, the $\pibar_g$, the column scaling and rank
reduction of $Z$, and hence $Z$, are fixed at each $n$; this uses that the cut points depend on
predictions only. $Z$ still changes with $n$ through the group means, so let $Q_n$ be a $G\times k_n$
matrix whose orthonormal columns span the columns of $Z$; then $S_n=\lVert Q_n^\top r\rVert^2$ and
\[
Q_n^\top r=\sum_{i=1}^n q_{g(i)}\,\frac{y_i-p_i}{\sqrt{V_{g(i)}}},
\]
where $q_g$ is the $g$th row of $Q_n$. Given the predictions this is a sum of independent, mean-zero
vectors whose covariance is exactly $Q_n^\top Q_n=I_{k_n}$, and each summand is bounded in norm by
$\max_g V_g^{-1/2}\to0$. For any fixed unit vector $a$, the scalar sums $a^\top Q_n^\top r$ therefore
satisfy Lindeberg's condition, and by the Cram\'er--Wold device $Q_n^\top r$ is asymptotically
$N(0,I_{k_n})$, uniformly over the rotation $Q_n$ because the bound does not involve it. Hence $S_n$ is
asymptotically $\chi^2_{k_n}$, in the uniform sense stated. Nothing in the argument uses equal group sizes
or a particular distribution of the $p_i$: every non-empty group need only keep growing. The condition
on the outcomes fails if the deployed model changes the care patients receive, or if who arrives
depends on earlier outcomes; cut points taken from the predictions of a first batch are allowed, since
they use predictions only. Under drift in the risk distribution the conclusion stands, although small groups weaken
the protection against corrupted records that equal-frequency groups give (Proposition~1 of the main
text bounds a record's influence by the variance of its groupmates).

(c) Finding the group of a new patient is a binary search among $G-1$ cut points, and updating four sums
is constant work. The statistic needs the $G$ residuals, the basis at $G$ points and a QR decomposition of
a $G\times k$ matrix. \hfill$\square$

\paragraph{Repeated looks.} Part (b) concerns one look. The statistics at successive looks share their
patients and are dependent, so acting on the first rejection over several looks does not have level
$\alpha$ at each look's $\chi^2_k$ critical value; splitting the level, $\alpha/L$ at each of $L$ looks
fixed in advance, keeps the overall false-alarm rate at most $\alpha$ in large samples; looks added
after the start break this bound. With fixed groups and stable group means the vector $r$ has
approximately independent increments between looks, so group-sequential boundaries for chi-squared
statistics \citep[chapter 15]{jennison2000group} apply and would be less conservative than splitting
the level.

\paragraph{Simulation.} The specification was fixed and hashed before any replicate was run. Reference
predictions were $10{,}000$ draws with $\mathrm{logit}\,p\sim N(-1.5,1)$, and their deciles gave the cut
points; patients then arrived in batches of $100$, $2000$ replicates a cell. Under a calibrated model the
level at $5\%$ was $0.047$ after $10{,}000$ patients and $0.054$ after $100{,}000$; when every monitored
patient had $\mathrm{logit}\,p\sim N(-0.5,1.3)$, a distribution different from the reference, so that the
smallest group held about $3900$ of $100{,}000$, it was $0.050$ and $0.052$. Testing after each of ten blocks of $1000$ patients and stopping at the first
rejection gave a false-alarm rate of $0.214$ at $5\%$ a look and $0.030$ at $0.5\%$ a look; the last
look alone read $0.057$. Against an overfitted model, $\mathrm{logit}\,\Pr(y=1)=0.8\,\mathrm{logit}\,p$,
the test rejected in $0.993$ of replicates after $2000$ patients and in all after $10{,}000$. The median time to add a hundred patients was about $0.5$
milliseconds and to compute the test about $1.5$ milliseconds, on one core of an AMD Ryzen~9 3900X; both
are independent of the number of patients already added, which was a million when they were timed.

\section{Power by departure family}\label{si:pairedfam}

Table~\ref{tab:pairedfam} replaces the counts of scenarios in which EDGE leads or ties a comparator by the
mean difference in size-adjusted power over each departure family, with a Monte Carlo interval. The counts
in Figures~1 and~2 of the main text use a margin of $0.01$, which is close to the Monte Carlo standard error
of a single difference, so they do not separate a small loss from a tie; the means do. The last column is
what ten groups cost against the default partition on clean data.

\input{tab_pairedfam}

\section{What protects a test}\label{si:tw}

\paragraph{Design.} The study was specified and hashed before any replicate ran, with two addenda written
before their replicates: one adding a milder diagnostics rule after a smoke run of three replicates showed
that the declared rule removes about $3.5\%$ of a clean data set, the other adding the tests of Hosmer and
Hjort. Every test runs on the same data sets, $1000$ replicates a cell, so all comparisons are paired. The
two-covariate design is that of the corruption study of the main text; the twelve-covariate design has
eight standard normal and four Bernoulli($0.3$) covariates at $n=2000$, with the error applied to the first
covariate. The ungrouped twin is the score test for adding the orthogonal cubic polynomial in the fitted
risk to the model, on three degrees of freedom. The tests of Hosmer and Hjort use their general-purpose
weight $\hat\pi\log\hat\pi$ at ten groups, in their Hosmer--Lemeshow form on $G-2$ degrees of freedom and
their full-covariance form; their implementation, written from the paper, held its level within
$0.05\pm0.02$ at $n=1000$ and $5000$ before the study.

\paragraph{Results.} Table~\ref{tab:tw} gives the rates. The grouped tests at ten groups, EDGE in its unit form and the Hosmer--Lemeshow statistic, stay below
$0.10$ in every setting with exaggerated covariates. Under sign errors EDGE at ten groups holds with one
reversed record and with ten among twelve covariates, and reads $0.210$ with ten reversed records in
$1000$ and $0.236$ with twenty-five in $5000$; the Hosmer--Lemeshow statistic reads $0.093$ and $0.120$
there. Every ungrouped test, whatever the scale of its directions, fails at one or two
corrupted records once the model is refitted. On frozen predictions (Table~\ref{tab:twext}) the ungrouped
twin holds against exaggerated covariates but not against five or more reversed records, and
Spiegelhalter's $z$ holds throughout. This frozen-prediction arm was added by a third addendum, hashed
before it ran, on the data sets of the external-validation study (S7), which it regenerates and checks
against that study's Spiegelhalter $p$-values.
Spiegelhalter's $z$ rejects in none of the $25{,}000$ replicates at its nominal level, because its normal
reference ignores the fitted coefficients; at its own realised level it is as fragile as Stukel's test. The
tests of Hosmer and Hjort group the records but weight each within its group by a function of its
covariates, and fail at about five reversed records in $1000$. EDGE's $p$-values from the
Satterthwaite and the Imhof reference gave identical rejections at $5\%$ in every cell.

\paragraph{Influence diagnostics.} Removing every record with Pregibon's $\Delta X^2>4$ or $\Delta\beta>1$,
the declared rule, removed $3.7\%$ of the records of a clean data set at $n=1000$ and every test then
rejected a correct model in more than $0.95$ of replicates, because the records removed are those whose
outcome is least likely under the model. Removing only the records with $\Delta\beta>1$, the addendum's
rule, removed none in any cell, so it changed nothing.

\paragraph{The Pulkstenis--Robinson test.} It was run in every scenario of the main battery and is reported
here. Its mean size-adjusted power over the $158$ alternatives is $0.16$, against $0.54$ for EDGE at the
default partition and $0.48$ for the Hosmer--Lemeshow statistic at ten groups. It groups on covariate
patterns \citep{pulkstenis2002two}; against exaggerated covariates it stays below $0.10$ up to ten records
in $1000$ and twenty-five in $5000$ and reaches $0.121$ at fifty in $5000$, and it is not protected against
sign errors: $0.224$ with ten reversed records in $1000$ and
$0.888$ with fifty in $5000$.

\input{tab_tw}

\input{tab_twext}

\section{Proof of Theorem 2}\label{si:proof2}

Notation and assumptions are those of Appendix A of the main text: $\varphi$ are the shape functions of
Theorem~2, $M=\int_0^1\varphi\varphi^\top dq$, $K=\int_0^1\varphi\,\kappa^\top dq$, and the limit to be
shown is (A1) of the main text.

\paragraph{Step 0: the bases} Because $P_Z$ depends on $Z$ only through its column space, and
(Step~6) so do the weights, the columns may be taken as $Z_{gj}=\varphi_j(q_g)$ once the bases of
Section~2.2 of the main text are shown to span that space up to a uniform $o(1)$. They evaluate their
shapes at $\pibar_g$ or $\bar\eta_g=\mathrm{logit}\,\pibar_g$ rather than at $q_g$. Under (A3) and
(A4$'$) the empirical quantile function of the fitted risk satisfies
$\sup_q|\hat Q_{\hat\pi}(q)-Q_\pi(q)|=O_p(n^{-1/2})$ and $Q_\pi$ is Lipschitz, so averaging over a
group of width $1/G$ gives $\pibar_g=Q_\pi(q_g)+O(1/G)+O_p(n^{-1/2})$ uniformly in $g$; with shape
functions continuously differentiable this gives $z_{gj}=\varphi_j(q_g)+\delta_{gj}$ with
$\sup_g\|\delta_g\|=O(1/G)+O_p(n^{-1/2})$, where $\varphi_j$ is the shape composed with $Q_\pi$ (or
with $Q_\eta$ for the Stukel columns). The contribution of $\delta$ to $\zeta_n$ below is a sum whose
summands are bounded by $Cn^{-1/2}\sup_g\|\delta_g\|$, handled as replacement (a) of Step~4, and its
contribution to $G^{-1}Z^\top Z$ and to Step~5 is $O(1/G)+O_p(n^{-1/2})$ by the same bound.
Orthonormalising the columns on the observed $\{\pibar_g\}$ multiplies $Z$ by an invertible matrix
that converges at the same rate, and $S$ and the weights are invariant to it. For the Stukel basis
linear independence needs $0$ interior to the support of $\eta$; otherwise one sign-split column
vanishes identically and is dropped.

\paragraph{Step 1: the statistic as a smooth-weighted sum} Put $\zeta_n=G^{-1/2}Z^\top r$, so that
$S=\zeta_n^\top(G^{-1}Z^\top Z)^{-1}\zeta_n$ with $G^{-1}Z^\top Z\to M$, a Riemann sum, and $M$
invertible because the $\varphi_j$ are linearly independent. By Step~1 of the proof of
Theorem~1 of the main text, $r=\tilde r-U(X^\top WX)^{-1}s+\rho$, hence
\begin{equation}\label{si:zeta}
\zeta_n=\sum_{i=1}^n w_i\varepsilon_i\;-\;\big[G^{-1/2}Z^\top U(X^\top WX)^{-1/2}\big]
\big[(X^\top WX)^{-1/2}s\big]\;+\;G^{-1/2}Z^\top\rho,
\qquad w_i=\frac{G^{-1/2}z_{g(i)}}{\sqrt{V_{g(i)}}}\in\mathbb{R}^d .
\end{equation}
Under (A3) the true linear predictor has compact support and $v_i\ge c_0>0$ for every $i$. Since
$\mathbb{E}\|x\|^2<\infty$ gives $\max_{i\le n}\|x_i\|=o_p(\sqrt n)$, and $\mathrm{expit}$ is
$\tfrac14$-Lipschitz, $\sup_i|\hat\pi_i-\pi_i|\le\tfrac14\max_i\|x_i\|\,\|\hat\beta-\beta\|=o_p(1)$,
so $\sup_i|\hat v_i/v_i-1|=o_p(1)$. On an event of probability tending to one, therefore,
$V_g\ge\tfrac12c_0|g|$ simultaneously in $g$, whatever the group sizes and including $|g|=1$, and
the normalisation computed from $\hat\pi$ differs from the same sum with $\pi$ by a factor
$1+o_p(1)$ uniformly in $g$. Since $\varphi$ is bounded and $|g|=n/G$,
$\|w_i\|\le C\,G^{-1/2}(n/G)^{-1/2}=C\,n^{-1/2}$ for every $i$, whatever the group size.

\paragraph{Step 2: the remainder} The second-order term of the expansion of $\hat\pi_i$ gives
$|\rho_g|\le C\,V_g^{-1/2}\sum_{i\in g}\|x_i\|^2\|\hat\beta-\beta\|^2$, so
$\|G^{-1/2}Z^\top\rho\|\le C\,n^{-1/2}\|\hat\beta-\beta\|^2\sum_i\|x_i\|^2=O_p(n^{-1/2})$. Only
$\sum_i\|x_i\|^2=O_p(n)$, which (A2) supplies, is used; no bound within a single group is needed.
The $o_p(n^{-1/2})$ of the expansion in (A2) contributes
$\|G^{-1/2}Z^\top U(X^\top WX)^{-1/2}\|\,O_p(\sqrt n)\,o_p(n^{-1/2})=o_p(1)$ by Step~5, and
replacing $v_i$ by the $\hat v_i$ carried in $U$ contributes at most
$\sum_i\|w_i\|\,|\hat v_i-v_i|\,\|x_i\|\,\|\hat\beta-\beta\|\le Cn^{-1/2}\,o_p(1)\,O_p(n)\,
O_p(n^{-1/2})=o_p(1)$ by Step~1.

\paragraph{Step 3: the central limit theorem with population weights} Let
$w_i^0=n^{-1/2}\varphi(q_i)/\sqrt{v_i}$ and $\xi_n^0=\sum_iw_i^0\varepsilon_i$. The pairs
$(w_i^0\varepsilon_i,\ n^{-1/2}J_0^{-1/2}x_i\varepsilon_i)$ are independent and identically
distributed, $w_i^0$ being a fixed function of $x_i$ through $q_i$, with finite variance by
$\mathbb{E}\|x\|^2<\infty$, so the Lindeberg--L\'evy theorem and the Cram\'er--Wold device give joint
asymptotic normality. Ranking by $\hat\pi$ is ranking by $\hat\eta$, and $F_\pi(\pi_i)=F_\eta(\eta_i)$,
so the quantile position of the predicted risk is $q_i$; by (A4$'$), $v_i=v(q_i)$. By the law of
large numbers $\mathrm{Var}(\xi_n^0)\to\mathbb{E}[\varphi(q)\varphi(q)^\top]=M$, $q$ being uniform,
and $\mathrm{Cov}(\xi_n^0,n^{-1/2}J_0^{-1/2}s)\to\mathbb{E}[\varphi(q)\sqrt{v(q)}\,x^\top]J_0^{-1/2}
=KJ_0^{-1/2}$. Hence $(\xi_n^0,\,n^{-1/2}J_0^{-1/2}s)\rightsquigarrow
N\big(0,\big(\begin{smallmatrix}M&KJ_0^{-1/2}\\ J_0^{-1/2}K^\top&I\end{smallmatrix}\big)\big)$.

\paragraph{Step 4: from population weights to the weights of the statistic} Three replacements
separate $\sum_iw_i\varepsilon_i$ from $\xi_n^0$: the position $q_i$ is replaced by the rank fraction
$R_i/n$; the rank fraction by its group position $q_{g(i)}$; and the ranks and the normalisation are
computed under $\hat\beta$ rather than $\beta$. (a)~Since $|\varphi(R_i/n)-\varphi(q_{g(i)})|\le C/G$,
at a fixed $b$ the second replacement changes the sum by a quantity with conditional mean zero and
conditional variance at most $C/G^2\to0$; its version uniform in $b$ is part of (c). (b)~At fixed $\beta$ the ranks are functions of $X$ alone, and
$\Delta_i=\varphi(R_i/n)-\varphi(q_i)$ satisfies $|\Delta_i|\le C\sup_t|\hat F_n(t)-F_\eta(t)|=
O_p(n^{-1/2})$ by the Dvoretzky--Kiefer--Wolfowitz inequality, so the conditional variance of
$\sum_in^{-1/2}\Delta_i\varepsilon_i/\sqrt{v_i}$ is $n^{-1}\sum_i\Delta_i^2=O_p(n^{-1})$. (c)~These
bounds hold at a fixed $b$, where the ranks are $X$-measurable; the statistic uses $\hat\beta$, so
uniformity over $\{\|b-\beta\|\le Cn^{-1/2}\}$ is needed, and $\hat\beta$ lies in that ball with
probability tending to one by (A2). Write $T_n(b)=\sum_i\{w_i(b)-w_i^0\}\varepsilon_i$. The weight
$w_i(b)$ depends on $b$ through the group index $g_b(i)$ and through $\hat\pi_i(b)$. The group
indices depend on $b$ only through the ordering of $\{x_i^\top b\}$, and the number of distinct
orderings is the number of cells of the arrangement of the $\binom n2$ hyperplanes
$\{b:x_i^\top b=x_j^\top b\}$ in $\mathbb{R}^{p}$, at most $\sum_{j\le p}\binom{\binom n2}{j}=
O(n^{2p})$ \citep{cover1965geometrical}: polynomially many, so a union bound replaces any chaining
argument. Within one cell, fixed at a point $b_0$, the group indices are constant and $b\mapsto
w_i(b)$ is $C^1$ through $\hat\pi_i(b)$, with first derivative bounded by $Cn^{-1/2}\|x_i\|$ and
second by $Cn^{-1/2}\|x_i\|^2$, so $T_n(b)-T_n(b_0)$ is a linear term in $b-b_0$ plus
$O_p(n^{-1/2})$. The coefficient of the linear term is $\sum_i\{\partial w_i(b_0)/\partial b\}\varepsilon_i$,
whose weights are $X$-measurable within the cell and bounded by $Cn^{-1/2}\|x_i\|$, so the union bound
below applies to it as well, uniformly over cells; multiplied by $\|b-b_0\|\le Cn^{-1/2}$ it is
$O_p(\sqrt{\log n}\,n^{-1/2})$. At $b_0$ both the linear term's coefficient and $T_n(b_0)$ are sums of
$X$-measurable weights times the centred $\varepsilon_i$, with
$\sum_i\{w_i(b_0)-w_i^0\}^2=O_p(G_n^{-2}+n^{-1})$. The two terms come from the two replacements: the
group position differs from the rank fraction by at most $1/(2G_n)$, which by (a) contributes
$O(G_n^{-2})$, and the rank fraction differs from the population position by at most
$\sup_t|\hat F_{n,b}(t)-F_b(t)|+C\|x_i\|\,\|b-\beta\|$, where the first term is $O_p(n^{-1/2})$ uniformly
in $b$ because the half-spaces $\{x^\top b\le t\}$ form a Vapnik--\v{C}ervonenkis class
\citep[Section 2.6]{vandervaartwellner1996}. The second term is $o_p(1)$ uniformly in $i$ under the
moment condition of (A2) but not $O_p(n^{-1/2})$ uniformly; the bound needs only its sum of squares,
$n^{-1}\sum_i\|x_i\|^2=O_p(1)$, which (A2) gives. Hoeffding's inequality conditionally on $X$ gives
sub-Gaussian tails for both at each $b_0$, and a union bound over the $O(n^{2p})$ cells gives
$\sup_{\|b-\beta\|\le Cn^{-1/2}}|T_n(b)|=O_p\{\sqrt{\log n}\,(G_n^{-1}+n^{-1/2})\}$, which is
$o_p(1)$ under the rate condition $G_n/\sqrt{\log n}\to\infty$ of Theorem~2 of the main text; the
same bound covers the normalisation $V_g$ under $\hat\beta$. The group indicators do not form a
class of fixed complexity as $G\to\infty$, but the ordering they induce does, which is the random-cell
problem of \citet{Moore1975} in a regime their fixed-cell argument does not cover. Together,
$\sum_iw_i\varepsilon_i=\xi_n^0+o_p(1)$.

\paragraph{Step 5: the fitting correction} The $g$-th summand of $Z^\top U$ is
$z_gV_g^{-1/2}\sum_{i\in g}\hat v_ix_i^\top$, so $G^{-1/2}Z^\top U(X^\top WX)^{-1/2}=\sum_iw_i\hat
v_ix_i^\top(X^\top WX)^{-1/2}$, a sum over all $n$ observations. Write $\psi=\varphi/\sqrt v$,
uniformly continuous on $[0,1]$. Then $w_i=n^{-1/2}\psi(q_{g(i)})\{1+o_p(1)\}$ uniformly in $i$,
because $\{V_g/|g|\}^{-1/2}=v(q_g)^{-1/2}+o_p(1)$ uniformly in $g$: every observation in $g$ has a
position within $1/G$ of $q_g$, $v\circ Q_\eta$ is uniformly continuous by (A3) and (A4$'$), and
$\sup_i|\hat v_i-v_i|=o_p(1)$. Likewise
$|q_{g(i)}-q_i|\le1/G+\sup_t|\hat F_n(t)-F_\eta(t)|+C\sup_i|\hat\eta_i-\eta_i|=o_p(1)$ uniformly in
$i$, so $\delta_n=\sup_i|\psi(q_{g(i)})-\psi(q_i)|=o_p(1)$ and
\[
\Big\|\sum_iw_i\hat v_ix_i^\top-n^{-1}\sum_i\varphi(q_i)\sqrt{v_i}\,x_i^\top\Big\|
\le\{\delta_n+o_p(1)\}\,n^{-1}\sum_i\|x_i\|=o_p(1).
\]
Since $n^{-1}\sum_i\varphi(q_i)\sqrt{v_i}x_i^\top\to K$ and $(X^\top WX)^{-1/2}=n^{-1/2}J_0^{-1/2}
\{1+o_p(1)\}$, the bracket in \eqref{si:zeta} converges to $KJ_0^{-1/2}$. The law of large numbers
used is across the whole sample and not within a group, which is why the number of observations per
group need not grow.

\paragraph{Step 6: conclusion} From \eqref{si:zeta} and Steps~2--5,
$\zeta_n=\xi_n^0-KJ_0^{-1/2}(n^{-1/2}J_0^{-1/2}s)+o_p(1)\rightsquigarrow N(0,\Sigma)$, and by the
continuous-mapping theorem $S\rightsquigarrow\zeta^\top M^{-1}\zeta$ with $\zeta\sim N(0,\Sigma)$,
which is $\sum_j\lambda_j\chi^2_{1,j}$ with $\lambda_j$ the eigenvalues of $M^{-1}\Sigma$. These do
not depend on the basis chosen for $\mathrm{col}(Z)$: for invertible $A$,
$\{(ZA)^\top ZA\}^{-1}(ZA)^\top\Omega ZA=A^{-1}\{(Z^\top Z)^{-1}Z^\top\Omega Z\}A$ is a similarity
transformation. The plug-in weights at $G_n$, the eigenvalues of
$(G^{-1}Z^\top Z)^{-1}(G^{-1}Z^\top\widehat\Omega Z)$, converge to the $\lambda_j$: $G^{-1}Z^\top Z\to
M$ as a midpoint Riemann sum, and $G^{-1}Z^\top\widehat\Omega Z\to M-KJ_0^{-1}K^\top$ by Step~5. The
bounds follow: $KJ_0^{-1}K^\top\succeq0$ gives $\Sigma\preceq M$ and hence $\lambda_j\le1$, while
$\Sigma\succeq0$ as the covariance of the limit of $\zeta_n$. A weight is zero exactly when
$c^\top\varphi(q)=\sqrt{v(q)}\{a_0+a_1Q_\eta(q)\}$ for some $a_0,a_1$, the population form of the
level shift and tilt the fit absorbs; a zero weight removes a degree of freedom from the limit and
changes neither $\sum_j\lambda_j$ nor $\sum_j\lambda_j^2$, so the Satterthwaite and Imhof evaluations
are unaffected. \qed

\bibliographystyle{plainnat}
\bibliography{refs}
\end{document}
